\chapter{序論}
\section{背景}
作物の草丈は生育段階、栄養状態、倒伏リスク、収量形成などを把握する重要な表現型である。定規等による地上計測は直接的である一方、多数個体を反復して測るには労力がかかる。UAV画像は広域を短時間で取得でき、画像処理や写真測量と組み合わせた高スループットな表現型計測を可能にする。UAV画像・LiDARを用いた作物高推定は、作物、撮影時期、地面モデル、視線・スキャン角度に依存することが報告されている\cite{cropheight,uav3dcrop}。
単眼深度推定は一枚のRGB画像から各画素の奥行きを推定するため、追加深度センサを必要としない。視覚基盤モデルの特徴を下流タスクへ転移することで、単眼深度推定にも適応が可能となった\cite{dinov2,dav2}。ただし単眼画像から尺度を一意に定めることは原理的に困難であり、学習分布や幾何的手がかりに依存する\cite{eigen2014}。屋外一般シーンで学習されたMetric Depthモデルは、葉群、細い茎、列状植生、露出土壌、遮蔽が混在する農作物画像で尺度誤差や局所平滑化を生じ得る。
\section{目的と貢献}
本研究はDepth Anything V2 SmallのMetric Depth出力を農作物UAV画像に適応し、絶対深度精度と局所深度変化がどのように推移するかを既実施3実験から検討する。貢献は、(1) Relative由来Metric head、未学習Metric Outdoor、1 epoch追加学習、独立した5 epoch runを共通validationの定義に注意して統合すること、(2)局所差分の振幅・符号・GT差分階級・画像端の詳細分析から全体局所MAEの限界を示すこと、(3) epoch 1--5曲線から絶対値と局所形状の非同値性、再現性、未検証事項を記録することである。本稿は実施済み実験の中間報告であり、稲の垂直草丈または独立test性能を報告するものではない。